\subsection{Application-level predictive evaluation on BreaKHis}
\label{subsec:breakhis-predictive}

We complement the controlled synthetic experiments with a real-image
classification experiment on the BreaKHis breast histopathology dataset.
The analysis is restricted to the 40$\times$ magnification in
order to avoid confounding the representation comparison with changes in
microscopic scale. The selected subset contains 1995 RGB images from
81 patients.

\input{table_dataset.tex}

Each image is converted to a scalar hematoxylin representation by H\&E
stain deconvolution. A Gaussian filter with fixed standard deviation
$\sigma=1$ pixels is applied, after which the image is
quantized into 5 levels by multi-Otsu thresholding. The
4-connected components of equal quantized value define the segmentation
used by all compared representations. Neither the segmentation procedure
nor any of its parameters uses the benign/malignant labels.

For every segmentation we compute the segmentation-size baseline, the
ordinary RAG, the regional contact multigraph, the HIG without cut
augmentation, the binary-incidence HIG, and the full HIG. The latter
includes the typed generator counts, normalized generator statistics,
regional-hole descriptors, incidence densities, integer incidence
multiplicities, and graph-level quantities described in
Section~\ref{subsec:hig-descriptors}. The cellular quantity
$\chi_{\mathrm{cell}}$ is excluded from the predictor because it is
identically equal to $2$.

All 1995 processed images satisfied the implementation-level
identity
\[
|V_2|-|V_1|+|V_0|=2.
\]
Across this real-image experiment, the intrinsic construction processed
a total of 18952116 image-region generators and
3527723 cut generators.

The prediction task is benign versus malignant histology. We use
regularized logistic regression with class-balanced training weights.
Continuous descriptors are standardized from the training data only.
The regularization parameter is selected by grouped inner
cross-validation. The outer evaluation uses
5-fold stratified group cross-validation repeated
5 times, with patient identity as the grouping variable;
consequently, images from one patient never occur simultaneously in
training and test sets.

Repeated out-of-fold probabilities are averaged per image before the
final metrics are computed. Uncertainty is estimated by
2000 patient-level bootstrap resamples.

\input{table_hig_class_statistics.tex}

\input{table_predictive_performance.tex}

The full HIG obtains a cross-fitted ROC--AUC of
0.717
(95\% CI [0.615, 0.810]).
For comparison, the RAG obtains
0.722
(95\% CI [0.644, 0.802]).
The paired full-HIG minus RAG difference is
-0.005
(95\% CI [-0.036,
+0.028]).

The ablation comparisons are reported in
Table~\ref{tab:paired-auc}. The contact-multigraph minus RAG difference
is +0.007
(95\% CI [-0.009,
+0.028]); the full-HIG minus
HIG-without-cuts difference is +0.001
(95\% CI [-0.012,
+0.014]); and the full-HIG minus
binary-incidence-HIG difference is -0.004
(95\% CI [-0.011,
+0.005]).
These paired comparisons are interpreted as representation-level
ablations rather than as tests of the intrinsic mathematical correctness
of the HIG.

\input{table_paired_auc.tex}